Deep neural networks build abstractions through a series of transformations across layers. However, the extent to which each individual layer modifies the representation remains poorly quantified, especially in pure Multi-Layer Perceptrons (MLPs). While metrics like Centered Kernel Analysis (\cka) exist to compare representations across different models, a focused study on the ``step size'' of representational change within a single model and how it scales with architectural constraints is missing. In this work, we investigate the representational similarity across layers in deep MLPs using linear \cka. We measure \transcap as the divergence ($1 - \text{\cka}$) between adjacent layers across various widths, datasets, and training states. Our experiments on \mnist and \cifar reveal that adjacent layers are highly similar (\cka $> 0.8$), confirming that individual layers perform incremental transformations. Furthermore, we find that wider layers possess higher \transcap, leading to greater divergence between layers. Finally, we show that training significantly increases this divergence compared to random initialization, indicating that networks learn to utilize their available \transcap to extract distinct features. Our findings suggest that layer redundancy is an inherent property of MLP architectures, constrained by their width, which has significant implications for model compression and architectural design.
